\begin{thebibliography}{99}
\footnotesize\itemsep=0pt

\bibitem{A88}
Aleksandrov A.G., A de {R}ahm complex of nonisolated singularities,
 \href{https://doi.org/10.1007/BF01077604}{\textit{Funct. Anal. Appl.}} \textbf{22} (1988), 131--133.

\bibitem{A}
Aleksandrov A.G., Nonisolated hypersurface singularities, in Theory of
 Singularities and its Applications, \textit{Adv. Soviet Math.}, Vol.~1, \href{https://doi.org/10.1070/sm1990v065n02abeh001164}{Amer.
 Math. Soc.}, Providence, RI, 1990, 211--246.

\bibitem{A05}
Aleksandrov A.G., Logarithmic differential forms, torsion differentials and
 residue, \href{https://doi.org/10.1080/02781070500128313}{\textit{Complex Var. Theory Appl.}} \textbf{50} (2005), 777--802.

\bibitem{AT}
Aleksandrov A.G., Tsikh A.K., Th\'eorie des r\'esidus de {L}eray et formes de
 {B}arlet sur une intersection compl\`ete singuli\`ere, \href{https://doi.org/10.1016/S0764-4442(01)02166-8}{\textit{C.~R.~Acad.
 Sci. Paris S\'er.~I Math.}} \textbf{333} (2001), 973--978.

\bibitem{B}
Barlet D., Le faisceau {$\omega'_{X}$} sur un espace analytique {$X$} de
 dimension pure, in Fonctions de plusieurs variables complexes,~{III}
 ({S}\'em. {F}rançois {N}orguet, 1975--1977), \textit{Lecture Notes in
 Math.}, Vol.~670, \href{https://doi.org/10.1007/BFb0064400}{Springer}, Berlin, 1978, 187--204.

\bibitem{BS}
Brasselet J.P., Sebastiani M., Brieskorn and the monodromy, \href{https://doi.org/10.5427/jsing.2018.18f}{\textit{J.~Singul.}}
 \textbf{18} (2018), 84--104.

\bibitem{Br}
Brieskorn E., Die {M}onodromie der isolierten {S}ingularit\"aten von
 {H}yperfl\"achen, \href{https://doi.org/10.1007/BF01155695}{\textit{Manuscripta Math.}} \textbf{2} (1970), 103--161.

\bibitem{Bru}
Brunella M., Some remarks on indices of holomorphic vector fields,
 \href{https://doi.org/10.5565/PUBLMAT_41297_17}{\textit{Publ. Mat.}} \textbf{41} (1997), 527--544.

\bibitem{CM1}
Corr\^ea M., da~Silva~Machado D., Residue formulas for logarithmic foliations
 and applications, \href{https://doi.org/10.1090/tran/7584}{\textit{Trans. Amer. Math. Soc.}} \textbf{371} (2019),
 6403--6420, \href{https://arxiv.org/abs/1611.01203}{arXiv:1611.01203}.

\bibitem{CM2}
Corr\^ea M., da~Silva~Machado D., G{SV}-index for holomorphic {P}faff systems,
 \href{https://doi.org/10.25537/dm.2020v25.1011-1027}{\textit{Doc. Math.}} \textbf{25} (2020), 1011--1027, \href{https://arxiv.org/abs/1611.09376}{arXiv:1611.09376}.

\bibitem{D}
Douai A., Tr\`es bonnes bases du r\'eseau de {B}rieskorn d'un polyn\^ome
 mod\'er\'e, \href{https://doi.org/10.24033/bsmf.2348}{\textit{Bull. Soc. Math. France}} \textbf{127} (1999), 255--287.

\bibitem{GS}
Granger M., Schulze M., Normal crossing properties of complex hypersurfaces via
 logarithmic residues, \href{https://doi.org/10.1112/S0010437X13007860}{\textit{Compos. Math.}} \textbf{150} (2014), 1607--1622,
 \href{https://arxiv.org/abs/1109.2612}{arXiv:1109.2612}.

\bibitem{G}
Greuel G.M., Der {G}auss--{M}anin-{Z}usammenhang isolierter {S}ingularit\"aten
 von vollst\"andigen {D}urchschnitten, \href{https://doi.org/10.1007/BF01352108}{\textit{Math. Ann.}} \textbf{214}
 (1975), 235--266.

\bibitem{Gu}
Guimaraes A.G., Polin\^omio de {B}ernstein--{S}ato de uma hipersuperficie com
 singularidade isolada, Ph.D.~Thesis, {I}CMC-{U}SP, S\~ao Carlos, 2002.

\bibitem{K83}
Kersken M., Der {R}esiduenkomplex in der lokalen algebraischen und analytischen
 {G}eometrie, \href{https://doi.org/10.1007/BF01455946}{\textit{Math. Ann.}} \textbf{265} (1983), 423--455.

\bibitem{K84}
Kersken M., Regul\"are {D}ifferentialformen, \href{https://doi.org/10.1007/BF01185193}{\textit{Manuscripta Math.}}
 \textbf{46} (1984), 1--25.

\bibitem{L}
L\^e D.T., Calcul du nombre de cycles \'evanouissants d'une hypersurface
 complexe, \href{https://doi.org/10.5802/aif.491}{\textit{Ann. Inst. Fourier (Grenoble)}} \textbf{23} (1973),
 261--270.

\bibitem{M}
Michler R., Torsion of differentials of hypersurfaces with isolated
 singularities, \href{https://doi.org/10.1016/0022-4049(94)00117-2}{\textit{J.~Pure Appl. Algebra}} \textbf{104} (1995), 81--88.

\bibitem{NT20b}
Nabeshima K., Tajima S., Testing zero-dimensionality of varieties at a point,
 \href{https://doi.org/10.1007/s11786-020-00484-y}{\textit{Math.Comput.Sci.}}, {t}o appear, \href{https://arxiv.org/abs/1903.12365}{arXiv:1903.12365}.

\bibitem{NT16a}
Nabeshima K., Tajima S., Computing {T}jurina stratifications of
 {$\mu$}-constant deformations via parametric local cohomology systems,
 \href{https://doi.org/10.1007/s00200-016-0289-4}{\textit{Appl. Algebra Engrg. Comm. Comput.}} \textbf{27} (2016), 451--467.

\bibitem{NT16b}
Nabeshima K., Tajima S., Solving extended ideal membership problems in rings of
 convergent power series via {G}r\"obner bases, in Mathematical Aspects of
 Computer and Information Sciences, \textit{Lecture Notes in Comput. Sci.},
 Vol.~9582, \href{https://doi.org/10.1007/978-3-319-32859-1_22}{Springer}, Cham, 2016, 252--267.

\bibitem{NT17a}
Nabeshima K., Tajima S., Algebraic local cohomology with parameters and
 parametric standard bases for zero-dimensional ideals, \href{https://doi.org/10.1016/j.jsc.2017.01.003}{\textit{J.~Symbolic
 Comput.}} \textbf{82} (2017), 91--122, \href{https://arxiv.org/abs/1508.06724}{arXiv:1508.06724}.

\bibitem{NT17b}
Nabeshima K., Tajima S., Computing {$\mu^*$}-sequences of hypersurface isolated
 singularities via parametric local cohomology systems, \href{https://doi.org/10.1007/s40306-016-0198-4}{\textit{Acta Math.
 Vietnam.}} \textbf{42} (2017), 279--288.

\bibitem{NT18}
Nabeshima K., Tajima S., Computation methods of logarithmic vector fields
 associated to semi-weighted homogeneous isolated hypersurface singularities,
 \href{https://doi.org/10.21099/tkbjm/1554170422}{\textit{Tsukuba~J. Math.}} \textbf{42} (2018), 191--231.

\bibitem{NT18b}
Nabeshima K., Tajima S., A new method for computing the limiting tangent space
 of an isolated hypersurface singularity via algebraic local cohomology, in
 Singularities in Generic Geometry, \textit{Adv. Stud. Pure Math.}, Vol.~78,
 \href{https://doi.org/10.2969/aspm/07810331}{Math. Soc. Japan}, Tokyo, 2018, 331--344.

\bibitem{NT19b}
Nabeshima K., Tajima S., Alternative algorithms for computing generic
 {$\mu^*$}-sequences and local {E}uler obstructions of isolated hypersurface
 singularities, \href{https://doi.org/10.1142/S0219498819501561}{\textit{J.~Algebra Appl.}} \textbf{18} (2019), 1950156,
 13~pages.

\bibitem{NT19a}
Nabeshima K., Tajima S., Computing logarithmic vector fields and
 {B}ruce--{R}oberts {M}ilnor numbers {\it via} local cohomology classes,
 \textit{Rev. Roumaine Math. Pures Appl.} \textbf{64} (2019), 523--540.

\bibitem{NT20}
Nabeshima K., Tajima S., Generalized integral dependence relations, in
 Mathematical Aspects of Computer and Information Sciences, \textit{Lecture
 Notes in Computer Science}, Vol.~11989, \href{https://doi.org/10.1007/978-3-030-43120-4_6}{Springer}, Cham, 2020, 48--63.

\bibitem{P}
Pol D., On the values of logarithmic residues along curves, \href{https://doi.org/10.5802/aif.3176}{\textit{Ann. Inst.
 Fourier (Grenoble)}} \textbf{68} (2018), 725--766, \href{https://arxiv.org/abs/1410.2126}{arXiv:1410.2126}.

\bibitem{P2}
Pol D., Characterizations of freeness for equidimensional subspaces,
 \href{https://doi.org/10.5427/jsing.2020.20a}{\textit{J.~Singul.}} \textbf{20} (2020), 1--30, \href{https://arxiv.org/abs/1512.06778}{arXiv:1512.06778}.

\bibitem{S71}
Saito K., Quasihomogene isolierte {S}ingularit\"aten von {H}yperfl\"achen,
 \href{https://doi.org/10.1007/BF01405360}{\textit{Invent. Math.}} \textbf{14} (1971), 123--142.

\bibitem{S73}
Saito K., Calcul alg\'ebrique de la monodromie, \textit{Ast\'erisque}
 \textbf{7} (1973), 195--211.

\bibitem{S77}
Saito K., On the uniformization of complements of discriminant loci, in
 Hyperfunctions and Linear Partial Differential Equations, \textit{RIMS
 K\^oky\^uroku}, Vol.~287, Res. Inst. Math. Sci. (RIMS), Kyoto, 1977,
 117--137.

\bibitem{S}
Saito K., Theory of logarithmic differential forms and logarithmic vector
 fields, \textit{J.~Fac. Sci. Univ. Tokyo Sect.~IA Math.} \textbf{27} (1980),
 265--291.

\bibitem{Sch}
Scherk J., On the {G}auss--{M}anin connection of an isolated hypersurface
 singularity, \href{https://doi.org/10.1007/BF01351450}{\textit{Math. Ann.}} \textbf{238} (1978), 23--32.

\bibitem{Sch93}
Scherk J., On the pole order and {H}odge filtrations of isolated hypersurface
 singularities, \href{https://doi.org/10.4153/CMB-1993-050-9}{\textit{Canad. Math. Bull.}} \textbf{36} (1993), 368--372.

\bibitem{Schu}
Schulze M., Algorithms for the {G}auss--{M}anin connection, \href{https://doi.org/10.1006/jsco.2001.0482}{\textit{J.~Symbolic
 Comput.}} \textbf{32} (2001), 549--564.

\bibitem{BSkoda}
Skoda H., Briançon J., Sur la cl\^oture int\'egrale d'un id\'eal de germes
 de fonctions holomorphes en un point de {${\bf C}^{n}$}, \textit{C.~R.~Acad.
 Sci. Paris S\'er.~A} \textbf{278} (1974), 949--951.

\bibitem{T}
Tajima S., On polar varieties, logarithmic vector fields and holonomic
 {D}-modules, in Recent Development of Micro-Local Analysis for the Theory of
 Asymptotic Analysis, \textit{RIMS K\^oky\^uroku Bessatsu}, Vol.~40, Res.
 Inst. Math. Sci. (RIMS), Kyoto, 2013, 41--51.

\bibitem{TN20}
Tajima S., Nabeshima K., An algorithm for computing torsion differential forms
 associated to an isolated hypersurface singularity, \href{https://doi.org/10.1007/s11786-020-00486-w}{\textit{Math.Comput.Sci.}},
 {t}o appear.

\bibitem{TNN}
Tajima S., Nakamura Y., Nabeshima K., Standard bases and algebraic local
 cohomology for zero dimensional ideals, in Singularities~--
 {N}iigata--{T}oyama 2007, \textit{Adv. Stud. Pure Math.}, Vol.~56, \href{https://doi.org/10.2969/aspm/05610341}{Math. Soc.
 Japan}, Tokyo, 2009, 341--361.

\bibitem{TN20a}
Tajima S., Shibuta T., Nabeshima K., Computing logarithmic vector fields along
 an {ICIS} germ via {M}atlis duality, in Computer Algebra in Scientific
 Computing, \textit{Lecture Notes in Computer Science}, Vol.~12291, \href{https://doi.org/10.1007/978-3-030-60026-6_32}{Springer},
 Cham, 2020, 543--562.

\bibitem{Te73}
Teissier B., Cycles \'evanescents, sections planes et conditions de {W}hitney,
 \textit{Ast\'erisque} \textbf{7} (1973), 285--362.

\bibitem{Te77}
Teissier B., Vari\'et\'es polaires. {I}.~{I}nvariants polaires des
 singularit\'es d'hypersurfaces, \href{https://doi.org/10.1007/BF01425742}{\textit{Invent. Math.}} \textbf{40} (1977),
 267--292.

\bibitem{vS}
van Straten D., The spectrum of hypersurface singularities,
 \href{https://arxiv.org/abs/2002.00519}{arXiv:2002.00519}.

\bibitem{V}
Vetter U., \"Au{\ss}ere {P}otenzen von {D}ifferentialmoduln reduzierter
 vollst\"andiger {D}urchschnitte, \href{https://doi.org/10.1007/BF01168480}{\textit{Manuscripta Math.}} \textbf{2}
 (1970), 67--75.

\bibitem{Z}
Zariski O., Characterization of plane algebroid curves whose module of
 differentials has maximum torsion, \href{https://doi.org/10.1073/pnas.56.3.781}{\textit{Proc. Nat. Acad. Sci. USA}}
 \textbf{56} (1966), 781--786.

\end{thebibliography}
